% RESEARCH_DRAFT_SHA256: 97d76efe556c13ffd2652adc295e219e38d9fd32172e9e7485ea460ae97df8a1
% Layout and source-reference wrappers expanded for the web renderer.
\documentclass[11pt]{article}
\usepackage[margin=25mm]{geometry}
\usepackage{fontspec,unicode-math}
\setmainfont{Latin Modern Roman}
\setmathfont{Latin Modern Math}
\usepackage{amsmath,amssymb,mathtools,xurl,hyperref}
\hypersetup{colorlinks=true,urlcolor=blue,linkcolor=blue}
\setcounter{secnumdepth}{0}
\setlength{\parindent}{0pt}
\setlength{\parskip}{5pt}
\setlength{\emergencystretch}{3em}
\newcommand{\sref}[2]{\hyperlink{source:#1:#2}{[S#2]}}
\newcommand{\cataloguescope}[1]{\paragraph{Source record.}#1}
\newcommand{\floor}[1]{\left\lfloor#1\right\rfloor}
\begin{document}
\section*{Research-attempt provenance}
This standalone research attempt retains the complete original definition body below,
including its historical source and preparation notes. Those notes describe the
earlier catalogue preparation; the new attempt and its verification are identified
separately after that body. No original ID, catalogue status, or source statement is changed.

\section{Certified landing-page staging for PreTeXt}\label{statement}

\subsection{Definitions and mathematical statement}

\paragraph{Shared notation.}Standard number systems are $\mathbb N,\mathbb Z,\mathbb Q,\mathbb R,\mathbb C$. Specialized notation and index ranges follow the cited source and the definitions below.


A landing page is an entry page that helps a visitor identify a resource and navigate to relevant content or actions. Navigation labels, audience needs, and the information architecture determine how it functions. A design evaluation requires explicit tasks or user goals.
\paragraph{Record scope and qualifications.}
The catalogue title is retained as metadata. Named certificates, contracts, bounds or frameworks in that title are not asserted as results of the original source extract.
The referenced page and the criterion for better are not identified in the short extract. Its description of existing support is historical.

Catalogue screening limitation: PreTeXt software design, documentation, or publication feature request; no mathematical problem is specified by the supplied passage.

The source snapshot labels this record partially solved; this is catalogue metadata rather than independent verification of a mathematical result or the current completion of a project.
\paragraph{Statement recorded in the supplied source.}
Landing page (this sort of exists, but it could be better). \href{https://github.com/PreTeXtBook/community/wiki/Moderated-Problem-Session}{[S2]}
\paragraph{Source status.}
The supplied catalogue labels this source record partially solved; that metadata was checked against the supplied snapshot on 2026-10-11. The label is preserved for provenance and is not an independent verification of present open status or mathematical correctness.
\subsection{Short English statement}
Improve the project's landing page.
\subsection{Sources}
\begin{itemize}
\item[S1] Ulam AI and the UnsolvedMath Contributors, UnsolvedMath: A Curated Collection of Open Mathematics Problems, catalogue entry. Catalogue curation and metadata: CC BY 4.0. The source extract is retained; definitions, English summaries and qualifications are editorial additions.
\url{https://huggingface.co/datasets/ulamai/UnsolvedMath}.
\item[S2] AIM Workshop Problem Lists, source linked by the catalogue.
\url{https://github.com/PreTeXtBook/community/wiki/Moderated-Problem-Session}.
\end{itemize}
\paragraph{Preparation and review.}
The original source and UnsolvedMath contributors retain attribution for the source extract. Codex (GPT-6; exact serving revision not exposed) assisted with definitions, English summaries, scope notes and formatting. No new human mathematical authorship or independently verified solution is asserted. Catalogue identity checked against the complete supplied registry and current repository. Distinct upstream records, including repeated questions and broad or damaged extracts, are retained for complete coverage of the source CS category. A repeated record is not a new independent conjecture. Complete mathematical-equivalence and current-literature review has not been performed. Preparation and final review follow contribution prompts C and D, with source gaps disclosed. Structural and rendering checks do not establish a new theorem or open status.
\label{end:statement}

\section{Literature and context check}
The community-session source was consulted directly. Its entry is a feature or design request, without a complete formal specification. The catalogue title is editorial and is not evidence of an existing theorem. The following finite model is an auxiliary investigation; no PreTeXt implementation or project completion is asserted.
\section{Research attempt: A landing-page update that exposes a mixed release}
Take a landing page $L$ linking to a target file $T$. Initially release 0 has $L_0$ linking to $T_0$; release 1 has $L_1$ linking to a new $T_1$. Publishing $L_1$ before making $T_1$ available creates a reachable broken link during that interval. Publishing the target first avoids this particular gap but does not ensure global consistency when several files refer to one another and readers fetch them at different times.

Here is a finite release model with a provable consistency property. Build every artifact for release $v$ under an immutable release-specific namespace. Verify that every internal link in its finite manifest resolves to a file in that namespace. Only after all files are available, atomically change one public pointer from $v_0$ to $v_1$. A reader resolves that pointer once and retains the chosen release identifier for all subsequent internal links. Assuming immutable available release files and atomic pointer reads, every link followed by that reader remains inside one complete release. Proof: the first fetch chooses one $v$; the manifest condition makes each internal edge end in an existing file of the same $v$; induction on the number of followed edges establishes the property.

Pinning only the landing page is insufficient if a child link is an unversioned latest URL. A reader who began with $v_0$ can then fetch a child from $v_1$, contradicting the same-release induction at that edge. The model exposes a precise staging requirement, but it does not select the landing page's navigation, audience, or accessibility design.
\section{Verification and adversarial audit}
The argument is an exact analysis of the stated finite model. Its counterexample was checked by direct substitution; no PreTeXt build, browser test, or user study was run. The model's assumptions are explicit and are not asserted to describe every real document.
\paragraph{Audit corrections.} Atomicity and release retention are assumptions of the abstract protocol, not claims about deployed hosting. Unversioned child links are tested as a counterexample to the proposed consistency proof.
\paragraph{External computation.} None; symbolic arguments and boundary checks only.
\section{Scope and remaining gaps}
The calculation identifies a concrete failed route or a restricted correctness property. Actual schema integration, transformations, usability, and compatibility with the full original feature request remain unimplemented and unverified. Thus the original agenda remains unresolved, with no claimed mathematical novelty.
\paragraph{Integration review.} Codex (GPT-6; exact serving revision not exposed) reviewed the written arguments and expanded layout and source-reference wrappers. Repository metadata and the production web renderer were checked using Python 3.13 and JavaScript with MathJax 3.2.2. This model-assisted review is not independent human proof verification; no PDF compilation or formal proof checking is asserted.

\section{Completion Estimate}
\noindent\textbf{COMPLETION: 0\%}
The calculations below the preserved source statement investigate a concrete route,
but establish no substantiated new progress toward the entire original target.
Known special cases and failed sufficient conditions are not credited as a solution.
\section{References}
The original source references and catalogue attribution are preserved in the
Sources section above. The following URLs were actually consulted for this attempt;
retention of an additional original reference does not claim fresh consultation of it.
\begin{itemize}
\item \url{https://github.com/PreTeXtBook/community/wiki/Moderated-Problem-Session} (consulted 2026-10-11). 
\end{itemize}
\section{Model and contributor disclosure}
OpenAI gpt-6-astra, requested reasoning configuration ultra, generated this research
attempt on 2026-10-11. The exact serving revision is not exposed. Codex GPT-6
performed orchestration, repository-format review, and a model-assisted mathematical
audit of the written arguments; its exact serving revision
is also not exposed. The mathematical route, boundary checks and final adversarial
self-audit followed contribution prompts A then B in
\texttt{docs/contribution\_prompts.md}, after reading \texttt{CONTRIBUTING.md}.
Python 3.14.2 assembled the standalone documents; any mathematical computation is
specified in Verification. No theorem prover was used. No human mathematical
authorship, substantial human proof correction, or independent human proof checking
is asserted. Model-assisted review, source retrieval, and structural checks do not
establish correctness or novelty. Full private reasoning transcripts are unavailable
as reproducibility artifacts; the complete arguments, inputs and outcomes actually
used are stated here or in the companion batch audit.
\section{Final status}
Unresolved as a new research contribution. This is a concrete route investigation,
not a claim of new complete resolution or a current-status certification. A known
answer reproduced for context does not become a new contribution. The catalogue's
source-status label is preserved; this attempt status does not assert that a
source-labelled solved question is open.
\end{document}
